%---------------------------------------------------------------------
% Brief edition: title page and a single orientation page.
%---------------------------------------------------------------------
\thispagestyle{empty}

% Depth measured from the rendered page, not computed: the compact edition's
% tighter leading still leaves the last masthead line below 10.4cm, where the
% amber rule would slice it in half.
\begin{tikzpicture}[remember picture, overlay]
  \fill[bmblued] (current page.north west) rectangle
        ([yshift=-11.3cm]current page.north east);
  \fill[bmaccent] ([yshift=-11.3cm]current page.north west) rectangle
        ([yshift=-11.5cm]current page.north east);
\end{tikzpicture}

\vspace*{0.9cm}
{\color{white}
  {\large علامہ اقبال اوپن یونیورسٹی، اسلام آباد}\par
  \vspace{3pt}
  {شعبۂ ریاضی \ \textbullet\ \ کلیۂ سائنس}\par
  \vspace{18pt}
  {\Huge\bfseries کاروباری ریاضی}\par
  \vspace{8pt}
  {\LARGE\bfseries اہم ترین سوالات --- مختصر ایڈیشن}\par
  \vspace{12pt}
  {\large کورس کوڈ \textenglish{1429} \ (بعض اوقات \textenglish{5405})}\par
  \vspace{4pt}
  % \par inside the group: colour is applied when the paragraph breaks, so
  % closing the group first would revert to black ink on a dark band.
  {نو یونٹ \ \textbullet\ \ سولہ سوالات \ \textbullet\ \ بی اے، بی کام، اے ڈی سی، بی ایس\par}
}

\vspace{1.1cm}

\begin{tcolorbox}[enhanced, colback=bmtint, colframe=bmblue, boxrule=0.6pt,
                  arc=3pt, left=12pt, right=12pt, top=10pt, bottom=10pt]
\textbf{\color{bmblue}یہ کتابچہ کیا ہے؟}\ \
اے آئی او یو کی سرکاری درسی کتاب \emph{کاروباری ریاضی} (کورس کوڈ
\textenglish{5405/1429}، اشاعت \textenglish{2023}) کے ہر یونٹ کے آخر میں دیے گئے
\textenglish{Self Assessment Questions} میں سے \textbf{سولہ سب سے اہم} سوالات کا
مختصر حل۔ یونیورسٹی اسائنمنٹ اور فائنل پرچہ اسی ذخیرے سے بناتی ہے۔

\medskip
\textbf{\color{bmblue}انتخاب کیسے ہوا؟}\ \
پرچے کے وزن کے مطابق۔ یونٹ \textenglish{6} (محدد، معکوس، کریمر) سے تین سوال؛
یونٹ \textenglish{1}، \textenglish{4}، \textenglish{7} اور \textenglish{9} سے دو
دو؛ باقی سے ایک ایک۔

\medskip
\textbf{\color{bmblue}تین ایڈیشن۔}\ \
یہ سب سے مختصر ہے۔ اس سے تفصیلی اردو ایڈیشن
(\textenglish{27} سوالات) \textenglish{AIOU\_1429\_Urdu\_Guess\_Paper.pdf} میں
اور مکمل انگریزی حل (\textenglish{84} صفحات)
\textenglish{AIOU\_1429\_Solved\_QuestionBank.pdf} میں ہیں۔

\medskip
\textbf{\color{bmblue}حیثیت۔}\ \
ذاتی مطالعے کے لیے۔ یہ اے آئی او یو کی مطبوعہ کتاب نہیں۔ اسے من و عن نقل کر کے
اسائنمنٹ میں جمع کرانا منع ہے۔
\end{tcolorbox}

\clearpage

%---------------------------------------------------------------------
\section*{پہلے یہ پڑھ لیجیے}
\markboth{پہلے یہ پڑھ لیجیے}{}

سوال پڑھیے۔ حل کو کاغذ سے ڈھانپیے۔ پہلے خود کاپی پر حل کیجیے، \emph{پھر}
ملائیے۔ صرف حل پڑھ لینے سے سیکھنے کا گمان ہوتا ہے؛ امتحان میں وہی کام آتا ہے جو
ناکام کوشش کے بعد سمجھ آیا ہو۔

\medskip
\begin{itemize}\setlength{\itemsep}{2pt}
  \item \textbf{سبز خانہ} --- آخری جواب، تاکہ طریقہ پڑھے بغیر اپنا نتیجہ ملا سکیں۔
  \item \textbf{پڑتال} --- جواب کو اصل مساوات میں واپس رکھ کر جانچ۔ اے آئی او یو
        طریقۂ کار کے نمبر دیتی ہے، اور پڑتال سے علامت کی غلطی ممتحن سے پہلے
        پکڑ میں آ جاتی ہے۔
  \item \textbf{سرخ خانہ} --- اس قسم کے سوال کی مخصوص غلطی۔ آخری صفحے پر ان کی
        مکمل فہرست ہے۔
\end{itemize}

\medskip
\noindent
\textbf{\color{bmblue}پرچے میں کس یونٹ کا کتنا وزن ہے}

\medskip
\begin{center}
\begin{tabular}{@{}p{4.4cm} p{1.8cm} p{1.6cm} p{6.6cm}@{}}
\toprule
\textbf{یونٹ اور موضوع} & \textbf{وزن} & \textbf{سوالات} & \textbf{اصل امتحان کس چیز کا ہے} \\
\midrule
\textenglish{1} احتمال            & زیادہ   & \textenglish{2} & جمع کا اصول بمقابلہ ضرب کا اصول \\
\textenglish{2} بے ترتیب متغیر    & درمیانہ & \textenglish{1} & احتمالات کا مجموعہ ہمیشہ \textenglish{1} \\
\textenglish{3} مساوات، عدم مساوات & درمیانہ & \textenglish{2} & علامت کا الٹنا؛ مطلق قدر کی دو صورتیں \\
\textenglish{4} خطی مساوات        & زیادہ   & \textenglish{2} & پہلے میلان، پھر ایک نقطہ \\
\textenglish{5} میٹرکس            & زیادہ   & \textenglish{1} & ضرب کے لیے ابعاد کا میل \\
\textenglish{6} محدد اور معکوس    & زیادہ   & \textenglish{3} & ہم عاملوں کی علامت؛ کریمر کا اصول \\
\textenglish{7} حد اور تفرق       & زیادہ   & \textenglish{2} & تقسیم کے اندر زنجیری اصول \\
\textenglish{8} جزوی مشتق         & درمیانہ & \textenglish{1} & دوسرے متغیر کو ثابت رکھنا \\
\textenglish{9} بہتر بندی         & زیادہ   & \textenglish{2} & دوسرے مشتق کی جانچ، ہمیشہ \\
\bottomrule
\end{tabular}
\end{center}

\medskip
\noindent
اگر وقت بہت کم ہو تو یونٹ \textenglish{6}، \textenglish{9}، \textenglish{4} اور
\textenglish{7} کو ترجیح دیجیے --- زیادہ تر نمبر اور تمام طویل سوالات انہی سے
آتے ہیں۔
